\section{习题 2.2：平面的方程（12 题）}\label{sec:22}
\tip{\textbf{基本工具：}有点 $P_0(x_0,y_0,z_0)$ 与法向量 $n=(A,B,C)$ 的平面为
\[ A(x-x_0)+B(y-y_0)+C(z-z_0)=0. \]
两方向向量 $u,v$ 不平行时，法向量 $n=u\times v$。已知三点则取两条差向量。一般仿射系中 $(A,B,C)$ 表示线性函数的系数，只有在直角正交坐标中才能无条件把它当作普通欧氏法向量。}
\begin{center}
\begin{tikzpicture}[scale=.84,>=Stealth]
\fill[blue!8] (0,0)--(5.2,.65)--(7.0,2.7)--(1.8,2.1)--cycle;
\draw[blue!65] (0,0)--(5.2,.65)--(7,2.7)--(1.8,2.1)--cycle;
\coordinate (P) at (3.2,1.1);\fill (P) circle(1.6pt) node[below left]{$P_0$};
\draw[->,very thick,blue] (P)--($(P)+(2.1,.2)$) node[right]{$u$};
\draw[->,very thick,teal] (P)--($(P)+(.8,1.25)$) node[above]{$v$};
\draw[->,very thick,red] (P)--($(P)+(-.5,2.1)$) node[above]{$n=u\times v$};
\node at (9.2,2) {$n\cdot (X-P_0)=0$};
\node at (9.2,1.3) {两个不共线方向确定平面};
\end{tikzpicture}
\end{center}

\Q{1}{哪些条件能唯一决定一张平面？}
(1) 过空间三个不同点：若三点共线，则有无穷多张平面通过该直线；应补充“三点不共线”。

(2) 过一点和一条直线：若点恰好在线上，则有无穷多张平面；应补充“点不在线上”。

(3) 过两条不同直线：若两线异面则根本不存在同时包含二者的平面；应补充“两直线共面”（相交或平行）。

(4) 过两点且平行于一条直线：若连接两点的方向与给定直线平行，则有无穷多张平面；应补充“两点连线与给定直线不平行”。
\ans{与书后答案 2.2 第 1 题一致。}

\Q{2}{给定点、方向向量或三点求平面方程}
\idea{找到两个独立方向 $u,v$，令 $n=u\times v$，再代入平面的点法式。以下各小问均可通过代入已知点和检查方向正交复验。}
\begin{enumerate}[label=(\arabic*)]
\item $P(1,1,1),u=(1,0,-1),v=(0,3,-4)$：
\[n=u\times v=(3,4,3),\quad 3(x-1)+4(y-1)+3(z-1)=0.\]
\ans{$3x+4y+3z-10=0$。}
\item $P_1(1,0,2),P_2(2,4,-1),P_3(-3,-5,1)$：
\[u=(1,4,-3),\ v=(-4,-5,-1),\quad u\times v=(-19,13,11).\]
将 $P_1$ 代入点法式并整体变号：
\ans{$19x-13y-11z+3=0$。}
\item 过 $(-1,0,2),(1,3,-2)$ 并平行 $u=(-1,-2,-4)$。两点差为 $v=(2,3,-4)$，
\[n=v\times u=(-20,12,-1).\]
取反写成 $n'=(20,-12,1)$，点法式得
\ans{$20x-12y+z+18=0$。}
\item 过 $(2,1,3)$ 并平行于 $2x-3y+z-4=0$：法向量不变 $n=(2,-3,1)$，常数项改为 $-4$，
\ans{$2x-3y+z-4=0$（此点恰好在原平面上，故两平面重合）。}
\end{enumerate}

\Q{3}{过原点、坐标轴或平行坐标轴的平面}
\begin{enumerate}[label=(\arabic*)]
\item 过 $O,( -2,1,2),(1,1,-2)$：两方向外积与 $(4,2,3)$ 平行，
\ans{$4x+2y+3z=0$。}
\item 过 $(1,2,-3)$，平行于 $xy$ 平面：$z$ 坐标恒为 $-3$，\ans{$z+3=0$。}
\item 过 $(1,2,0)$ 与 $y$ 轴：平面包含 $(0,1,0)$ 方向与 $(1,0,0)$ 方向，\ans{$z=0$。}
\item 平行 $x$ 轴，过 $(1,-1,2),(2,0,-1)$：线段方向 $(1,1,-3)$ 与 $e_x$ 的外积对应 $(0,3,1)$，\ans{$3y+z+1=0$。}
\item 经过 $z$ 轴、平行 $(1,-1,3)$：由两方向 $(0,0,1),(1,-1,3)$ 可得法向量 $(1,1,0)$，\ans{$x+y=0$。}
\end{enumerate}

\Q{4}{三点确定平面的四阶行列式表达式}
设 $P_i=(x_i,y_i,z_i)$（$i=1,2,3$）不共线。给出的行列式
\[
\begin{vmatrix}x&y&z&1\\x_1&y_1&z_1&1\\x_2&y_2&z_2&1\\x_3&y_3&z_3&1\end{vmatrix}=0
\]
第一行减去第二行，再展开得到\(\det[P-P_1;P_2-P_1;P_3-P_1]=0\)。它等价于 $P-P_1$ 与两独立边向量共面，所以恰是三点所确定的平面方程。
\ans{得证；“三点不共线”用于保证行列式不是恒等于 $0$。}

\Q{5}{积为零、平方差为零的图形}
令 $F=x+3y-z-2,\ G=3x+y+z-4$。
\begin{enumerate}[label=(\arabic*)]
\item $FG=0$ 等价于 $F=0$ 或 $G=0$；法向量 $(1,3,-1)$ 与 $(3,1,1)$ 不成比例，因此是\textbf{两张相交平面}的并。
\item 设 $F=2x+3y+z-1,G=-x+y-z+4$，则 $F^2-G^2=(F-G)(F+G)=0$；两个因子仍是非平行的一次式，所以也是\textbf{两张相交平面}的并。
\end{enumerate}
\ans{两小问都为两张相交平面；并集与两式同时为零的交集不能混淆。}

\Q{6}{判别两张平面平行、相交或重合}
看法向量是否成比例；若成比例，再看常数项是否同倍数。
\begin{enumerate}[label=(\arabic*)]
\item $(1,3,-1,-2)$ 与 $(2,6,-2,-2)$：法向量恰好二倍，常数项却不是二倍，故\textbf{平行但不重合}。
\item $x+y+3z-4=0$ 与 $x+3y+z-4=0$：法向量不成比例，故\textbf{相交}。
\item $3x+2y+z-6=0$ 与 $\frac x2+\frac y3+\frac z6-1=0$：后一式乘 $6$ 就是前式，故\textbf{重合}。
\end{enumerate}
\ans{依次为平行、相交、重合，与书后答案一致。}

\Q{7}{三张平面能否交于唯一点}
\idea{三张平面组成 $N\,X=-d$，唯一交点的充要条件是三阶系数矩阵 $N$ 可逆，即 $\det N\ne0$。}
(1) 三个法向量 $(1,1,-2),(1,-2,1),(-2,1,1)$ 的和为零，故行列式为零，无唯一交点。

(2) 前两张平面 $x+y+3z+4=0$ 与 $x+y+3z-4=0$ 平行且不重合，当然无共同点。

(3) 系数行列式 $\det\begin{pmatrix}1&-2&3\\1&2&1\\1&1&-1\end{pmatrix}\ne0$，故唯一交于一点。
\ans{(1) 否，(2) 否，(3) 是；与书后答案一致。}

\Q{8}{三张平面两两相交却不共点}
\[
\pi_1:ax+y+z+1=0,\quad \pi_2:x+ay+z+2=0,\quad\pi_3:x+y-2z+3=0.
\]
三法向量构成矩阵
\[
N=\begin{pmatrix}a&1&1\\1&a&1\\1&1&-2\end{pmatrix},\qquad\det N=-2(a-1)(a+2).
\]
若三张平面不能有唯一交点，必有 $a=1$ 或 $a=-2$。当 $a=1$ 时，前两张平面平行但不重合，不符合“两两相交”。当 $a=-2$ 时，三法向量两两不成比例，但三方程不相容，符合题意。
\ans{$a=-2$，与书后答案一致。}

\Q{9}{平面线性式与两点连线平行的关系}
设 $f(X)=Ax+By+Cz+D$，则
\[
f(M_1)-f(M_2)=A(x_1-x_2)+B(y_1-y_2)+C(z_1-z_2).
\]
右侧是平面系数线性式对方向向量 $M_1M_2$ 的作用，恰为零当且仅当该方向平行于给定平面；故所求等价关系成立。
\ans{对仿射坐标也成立；证明只需要线性方程的定义。}

\Q{10}{三张平行平面的交点有向分比}
三张平行平面写为 $f(X)+D_i=0$，其中 $f(X)=Ax+By+Cz$。直线参数化为 $X=X_0+t u$，代入得
\[t_i=-\frac{f(X_0)+D_i}{f(u)},\qquad t_P-t_Q=\frac{D_2-D_1}{f(u)}.\]
对分比 $(P,Q,R)=(t_R-t_P)/(t_Q-t_R)$ 的原书定义（注意端点排列），逐项代入并约去斜率，得到
\ans{$\displaystyle(P,Q,R)=\frac{D_1-D_3}{D_3-D_2}$，与书后答案一致。}

\Q{11}{到两张平行平面距离相等的点的轨迹}
两张平行平面同系数写为 $Ax+By+Cz+D_i=0$。对于同一点 $X$，两边距离分别正比于 $|f(X)+D_i|$，于是
\[
|f(X)+D_1|=|f(X)+D_2|.
\]
当 $D_1\ne D_2$ 时，移项平方后只有中间平面
\ans{$\displaystyle Ax+By+Cz+\frac{D_1+D_2}{2}=0$。}
\tip{与第一章的分点类似：到两平面等距，就是处于两条等值层之间的中点层。}

\Q{12}{到两平行平面的距离之比为 2:1}
给定 $\pi_1:3x-2y+5z+2=0$，$\pi_2:6x-4y+10z-5=0$。统一系数令 $L=6x-4y+10z$，第一式变成 $L+4=0$，第二式是 $L-5=0$。距离比应写为
\[
\frac{|L+4|}{|L-5|}=2\quad\Longleftrightarrow\quad L+4=\pm2(L-5).
\]
分别得到 $L=14$ 或 $L=2$。
\ans{两张轨迹平面分别为 $6x-4y+10z-14=0$、$6x-4y+10z-2=0$；与书后答案一致。}
